%%%%%%%%%%%%%%%%%%%%%%%%%%%%%%%%%%%%%%%%%%%%%%%%%%%%%%%%%%%%%%%%%%%%%%%%%%%%%%%%
% BLOQUE 3 — INTRODUCCIÓN Y CONTEXTO
%%%%%%%%%%%%%%%%%%%%%%%%%%%%%%%%%%%%%%%%%%%%%%%%%%%%%%%%%%%%%%%%%%%%%%%%%%%%%%%%

\phantomsection
\addcontentsline{toc}{section}{Introducción y contexto}
\section*{Introducción y contexto}

% TODO (equipo): presentar el tema TI-05 (serverless y edge computing),
% la empresa asignada (Terabyte), el alcance del trabajo y la
% estructura del informe. Las definiciones conceptuales de base van en
% el Marco teórico (bloque 4), no aquí.
%
% PÁRRAFO OBLIGATORIO DE CIERRE (punto 4 de las indicaciones — NO
% OMITIR): un párrafo breve, al cierre de esta sección, que distinga
% explícitamente qué contenido proviene directamente de la ficha TI-05
% y cuál es aporte propio del grupo (subtemas agregados, alternativas
% adicionales investigadas, mediciones propias, etc.). Este párrafo se
% usa como referencia al evaluar el análisis crítico del grupo.
% Autoría: humana obligatoria (punto 6.1).

\clearpage
